% ============================================================================
%  Public CV — Muhammad Abdullah Zahid
%  Build (inside cv/; latexmk fails under MiKTeX here, so run pdflatex twice):
%          pdflatex -output-directory=build cv.tex   (twice)
%          then copy build/cv.pdf to ../assets/cv.pdf
%  This is the version linked from the website. It mirrors the full CV but
%  intentionally leaves out home address, phone number, date of birth, GPA,
%  scholarships, and honors and awards; those stay in the private CV.
% ============================================================================
\documentclass[11pt,a4paper]{article}

\usepackage[T1]{fontenc}
\usepackage[utf8]{inputenc}
\usepackage[charter]{mathdesign}
\usepackage[a4paper,top=2cm,bottom=2cm,left=2.1cm,right=2.1cm]{geometry}
\usepackage[dvipsnames]{xcolor}
\usepackage{titlesec}
\usepackage{enumitem}
\usepackage{microtype}
\usepackage{fancyhdr}
\usepackage{needspace}
\usepackage[hidelinks]{hyperref}

% Bullets may break across pages, but never leave a single line stranded at either end.
\makeatletter\@clubpenalty=10000\makeatother
\clubpenalty=10000
\widowpenalty=10000
\raggedbottom
% Never leave a section heading alone at the bottom of a page.
\let\cvorigsection\section
\renewcommand{\section}[1]{\Needspace{6\baselineskip}\cvorigsection{#1}}

\definecolor{cvred}{RGB}{122,0,0}
\definecolor{cvlink}{RGB}{190,20,70}
\hypersetup{colorlinks=true,urlcolor=cvlink,linkcolor=cvlink,
  pdftitle={CV - Muhammad Abdullah Zahid},pdfauthor={Muhammad Abdullah Zahid}}

\setlength{\parindent}{0pt}
\setlength{\parskip}{0pt}

% Section headings: dark-red small caps with a rule underneath.
\titleformat{\section}{\color{cvred}\scshape\large}{}{0em}{}[{\color{cvred}\titlerule[0.6pt]}]
\titlespacing*{\section}{0pt}{14pt}{8pt}
\titleformat{\subsection}{\color{cvred}\normalsize}{}{0em}{}
\titlespacing*{\subsection}{0pt}{8pt}{3pt}

\setlist[itemize]{leftmargin=1.4em,label=\textbullet,itemsep=1pt,topsep=2pt,parsep=0pt,before={\interlinepenalty0\relax}}

\pagestyle{fancy}
\fancyhf{}
\renewcommand{\headrulewidth}{0pt}
\fancyfoot[R]{\footnotesize\color{gray}\thepage}
\fancyfoot[L]{\footnotesize\color{gray}Last updated November 2026}

% \entry{Organisation}{, Role}{Dates}
\newcommand{\entry}[3]{\par\Needspace{4\baselineskip}\textbf{#1}#2\hfill #3\par\nopagebreak}
% \sub{italic line under an entry}
\newcommand{\sub}[1]{{\itshape #1}\par}
% \pub{label}{title}{authors}{venue and status}
\newcommand{\pub}[4]{\par\vspace{5pt}\Needspace{3\baselineskip}{\bfseries[#1] #2}\par #3\par{\itshape #4}\par}
% \me — own name in the author list
\newcommand{\me}{\textbf{M.\,A.\ Zahid}}
% \item-style labelled line: \lab{Label} text
\newcommand{\lab}[1]{\textbf{#1}\ }

\begin{document}

% --- Header -----------------------------------------------------------------
\begin{center}
  {\LARGE\bfseries Muhammad Abdullah Zahid}\\[4pt]
  Undergraduate Researcher, NISAR Laboratory, Kyoto University of Advanced Science\\[3pt]
  \href{mailto:2023m654@kuas.ac.jp}{2023m654@kuas.ac.jp} (university)\quad$|$\quad
  \href{mailto:abdullahzahid1440@gmail.com}{abdullahzahid1440@gmail.com} (personal)\\[2pt]
  \href{https://zahidmabdullah.github.io}{zahidmabdullah.github.io}\quad$|$\quad
  \href{https://github.com/ZahidMAbdullah}{github.com/ZahidMAbdullah}\quad$|$\quad
  \href{https://www.linkedin.com/in/muhammad-abdullah-zahid-3b5a1226a/}{LinkedIn}
\end{center}

% --- Research interests -------------------------------------------------------
\section{Research Interests}

\lab{Mixed Reality Interfaces for Fine Robotic Telemanipulation.} How the operator's visual
representation of the robot and its workspace shapes fine manipulation, and whether spatially registered
mixed reality helps operators not only perform better but also acquire skills faster.

\vspace{4pt}
\lab{Skill Transfer from Virtual Training to Physical Execution.} Whether a single interface shared by
virtual training and real teleoperation lets skills carry over without readaptation, and how digital
twins can keep that interface faithful to a changing remote workspace in real time.

\vspace{4pt}
\lab{Wearable Haptics for Telemanipulation and XR.} Which manipulation events are worth conveying
through touch, and how lightweight wearables, vibrotactile or passive kinesthetic, can render them
without tethering the operator to a grounded device.

\vspace{4pt}
\lab{Autonomous Robot Control.} Robust semantic object navigation for mobile robots in dynamic,
uncertain environments.

% --- Education ----------------------------------------------------------------
\section{Education}

\entry{Kyoto University of Advanced Science (KUAS)}{, Kyoto, Japan}{Sep.\ 2027 -- Sep.\ 2029}
\sub{Master of Engineering \quad$|$\quad Admitted; enrolling Sep.\ 2027 (two-year program)}
\sub{Graduate School of Engineering \quad$|$\quad Division of Mechanical and Electrical Systems Engineering}
\begin{itemize}
  \item \lab{Research Supervisor} \href{https://www.kuas.ac.jp/en/f-r/eng/sajid-nisar/}{Dr.\ Sajid Nisar},
        \href{https://nisarlab.org/}{Novel Intelligent Systems \& Advanced Robotics Laboratory}
  \item \lab{Planned Research} Bridging Training and Execution: A Unified Mixed Reality Interface for
        Robotic Telemanipulation
\end{itemize}

\vspace{3pt}
\entry{Kyoto University of Advanced Science (KUAS)}{, Kyoto, Japan}{Sep.\ 2023 -- Sep.\ 2027}
\sub{Bachelor of Engineering \quad$|$\quad Expected graduation: Sep.\ 2027}
\sub{Faculty of Engineering \quad$|$\quad Department of Mechanical and Electrical Systems Engineering}
\begin{itemize}
  \item \lab{Research Supervisor} \href{https://www.kuas.ac.jp/en/f-r/eng/sajid-nisar/}{Dr.\ Sajid Nisar},
        \href{https://nisarlab.org/}{Novel Intelligent Systems \& Advanced Robotics Laboratory}
        (Mar.\ 2025 -- Present)
\end{itemize}

% --- Research experience --------------------------------------------------------
\section{Research Experience}

\entry{NISAR Laboratory, KUAS}{, Undergraduate Researcher}{Mar.\ 2025 -- Present}
\sub{Research Supervisor: Dr.\ Sajid Nisar \quad$|$\quad Novel Intelligent Systems \& Advanced Robotics Laboratory}
\begin{itemize}
  \item \lab{HaptLinX: Haptic-Enabled MR Telemanipulation Training (2025) [C4].} First prototype: a
        digital twin of a 6-DoF Kinova Gen3 Lite, controlled in MR on a Meta Quest 3S through hand
        tracking and mirrored in real time by the physical arm, with bHaptics glove and suit cues for
        contact and alignment. Selected for the IEEE World Haptics 2025 Student Innovation Challenge.
  \item \lab{MR Training Interface for Fine Telemanipulation (Lab Project, 2026 -- Present) [C1].}
        Built the training module of the unified MR interface and evaluated it in a twelve-participant
        $2\times2$ within-subject study against a 2D multi-view display under the same control mapping.
        MR raised performance (5.19 vs.\ 1.19 successful placements/trial) and learning rate
        ($+0.479$ placements/trial); vibrotactile cues cut self-collisions by 0.71 per trial but did not
        accelerate learning.
  \item \lab{Unified MR Interface from Training to Execution (In Progress; Master's Research from Sep.\ 2027).}
        Extending the same MR interface and hand-based control from virtual training to teleoperation of
        the physical robot through its digital twin, with the robot workspace synchronised through RGB-D
        object reconstruction and pose tracking (NVIDIA FoundationPose), to measure how skills transfer
        to real execution.
  \item \lab{Escapement-Like Kinesthetic Rendering for a Virtual Piano [C2].} Designed a
        displacement-dependent controller that renders a piano key's let-off through a passive wearable
        tendon-resistance device with no key surface. Force-proxy measurements over five device fittings
        showed a consistent drop at let-off, then recovery; ten participants detected the release with
        95\% median accuracy and located it within the keystroke with 80\%.
  \item \lab{Wearable vs.\ World-Grounded Haptics for Space Teleoperation [C3].} Contributed to a VR
        comparison of a world-grounded (sigma.7) and a wearable (Dexmo) interface, with and without force
        feedback, on a fine-manipulation task modelled on space robotics.
  \item \lab{Uncertainty-Aware Semantic Object Navigation (2026 -- Present).} Developing an ObjectNav
        framework that combines object memory, detection confidence, and spatial uncertainty with active
        exploration, in NVIDIA Isaac Sim.
\end{itemize}

% --- Publications ---------------------------------------------------------------
\section{Publications}

\subsection{Peer-Reviewed Conference Papers}
\pub{C1}{Skill Acquisition in Robotic Teleoperation: Effects of Mixed Reality and Vibrotactile Feedback}
    {\me, R.\,M.\ Suarez Flores, S.\ Nisar}
    {2026 9th International Conference on Robotic Systems and Applications (ICRSA 2026). IEEE Xplore indexing in progress.}
\pub{C2}{Escapement-Like Kinesthetic Rendering with a Passive Wearable Device for a Virtual Piano}
    {\me, S.\ Shukla, R.\,M.\ Suarez Flores, S.\ Nisar}
    {2026 9th International Conference on Robotic Systems and Applications (ICRSA 2026). IEEE Xplore indexing in progress.}
\pub{C3}{Wearable vs.\ World-Grounded Haptics for Space Teleoperation: A VR-Based Comparative Analysis of Kinesthetic Feedback for Fine Telemanipulation}
    {K.\,H.\ Joshi, R.\,M.\ Suarez Flores, \me, S.\ Nisar}
    {31st International Symposium on Artificial Life and Robotics 2026 (AROB 2026).}

\subsection{Non-Peer-Reviewed Papers}
\pub{C4}{Developing a Haptic-Enabled Mixed-Reality System for Fine Robotic Telemanipulation Training}
    {\me, R.\,M.\ Suarez Flores, K.\ K C, K.\ Meitani, A.\ Majumder, S.\ Nisar}
    {43rd Annual Conference of the Robotics Society of Japan (RSJ2025).}

% --- Presentations ----------------------------------------------------------------
\section{Presentations}

\lab{Oral} ICRSA 2026, Huizhou, China --- [C1] and [C2] (Oct.\ 30 -- Nov.\ 1, 2026).\\
\lab{Oral} RSJ2025, International Session --- Human and Robot, Tokyo, Japan --- [C4] (Sep.\ 2025).\\
\lab{Poster} IEEE World Haptics Conference 2025, Student Innovation Challenge --- Mixed Reality, Suwon,
South Korea --- \textit{HaptLinX} (Jul.\ 2025).

% --- Teaching -----------------------------------------------------------------------
\section{Teaching Experience}

\entry{Introduction to Robotic Manipulators}{, Teaching Assistant, KUAS}{Sep.\ 2026 -- Present}
\entry{Algorithmic Thinking and Programming with Python}{, Teaching Assistant, KUAS}{Apr.\ -- Jul.\ 2026}

% --- Projects -----------------------------------------------------------------------
\section{Projects and Leadership}

\entry{LLM-Driven Natural Language Robot Control}{, Team Leader, with AsiaQuest Co., Ltd.}{2025}
\begin{itemize}
  \item Led a five-member team building a mobile robot that executes spoken multi-step commands locally
        (Whisper, Mistral 7B, Grounded SAM 2 on RGB-D, ROS 2 SLAM and Nav2); worked on the language
        processing, vision, and navigation pipelines.
\end{itemize}

\entry{Kyorabo}{, Co-Founder}{2026 -- Present}
\begin{itemize}
  \item Research-collaboration matching app built by a team of five (React Native, Supabase), live on
        iOS and the web; reached 80 researchers from 45 institutions in 18 countries in its first week.
        \href{https://kyorabo.com/}{[Website]}
        \href{https://devpost.com/software/kyorabo-find-research-collaborators}{[Devpost]}
\end{itemize}

\entry{AirKeys}{, KUAS Cornerstone Project}{2025 -- 2026}
\begin{itemize}
  \item Self-initiated virtual piano: a VR version with hand tracking (2025), extended to MR with AI
        features and haptic feedback (2026); the origin of the haptic rendering work in [C2].
\end{itemize}

\entry{Mobile Robots}{, Mechatronics Laboratory}{2025 -- 2026}
\begin{itemize}
  \item \textit{ELIBot}: four-wheeled obstacle-course robot teleoperated through ROS 2 on a Raspberry
        Pi 4, with QR-code and traffic-light detection and STM32 motor control
        \href{https://github.com/ZahidMAbdullah/ELIBot}{[GitHub]}; \textit{self-balancing line follower}
        with staged LQR control
        \href{https://github.com/ZahidMAbdullah/Staged-LQR-LineFollowerBot}{[GitHub]}.
\end{itemize}

% --- Memberships & service -----------------------------------------------------------
\section{Memberships and Service}

\lab{Memberships} IEEE, Student Member; Robotics Society of Japan, Student Member (2025 -- Present).\\
\lab{Outreach} Featured Speaker, ``Study in Japan: Fully Funded Undergraduate Scholarships,'' Pakistan
Student Association Japan (PSAJ), online (Mar.\ 2026).

% --- Skills -----------------------------------------------------------------------------
\section{Skills}

\lab{XR and Haptics} Unity (C\#), Meta Quest 3S, digital twins, hand tracking, inverse kinematics,
bHaptics TactSuit Pro / TactGlove DK2, sigma.7, Dexmo.\\
\lab{Robotics} Kinova Gen3 Lite, teleoperation, ROS 2, NVIDIA Isaac Sim, computer vision, mobile robot
navigation, embedded control (Arduino, STM32).\\
\lab{Research Methods} Within-subject experiment design, user studies, statistical analysis, \LaTeX.\\
\lab{Programming} Python, C/C++; Linux, Git.\\
\lab{Applied AI} Integrating pretrained models: Grounded SAM 2, SAM 3D, FoundationPose, Whisper,
Mistral 7B.\\
\lab{Languages} English (IELTS 7.5), Japanese (JLPT N4).\\
\lab{Certifications} JSPS Research Ethics, eL CoRE (2026)
\href{https://zahidmabdullah.github.io/assets/certificates/zahid-2026-jsps-elcore-research-ethics-certificate.pdf}{[Certificate]};
MITx 6.00.1x, Computer Science and Python (2024)
\href{https://courses.edx.org/certificates/9ee1f4959cd14ef0beca032932c73867}{[Certificate]}.

\end{document}
